\documentclass[10pt,a4paper]{amsart}
\usepackage{amsmath,amssymb,amsthm,graphicx}
\usepackage{verbatim}
\newenvironment{reference}{\comment}{\endcomment}
%\newenvironment{reference}{}{}
\newenvironment{slogan}{\comment}{\endcomment}
\newenvironment{history}{\comment}{\endcomment}

% For commutative diagrams we use Xy-pic
\usepackage[all]{xy}

% We use 2cell for 2-commutative diagrams.
\xyoption{2cell}
% Theorem environments.
%
\theoremstyle{plain}
\newtheorem{theorem}[subsection]{Theorem}
\newtheorem{proposition}[subsection]{Proposition}
\newtheorem{lemma}[subsection]{Lemma}

\theoremstyle{definition}
\newtheorem{definition}[subsection]{Definition}
\newtheorem{example}[subsection]{Example}
\newtheorem{exercise}[subsection]{Exercise}
\newtheorem{situation}[subsection]{Situation}

\theoremstyle{remark}
\newtheorem{remark}[subsection]{Remark}
\newtheorem{remarks}[subsection]{Remarks}

\numberwithin{equation}{subsection}

% Macros
%
\def\lim{\mathop{\mathrm{lim}}\nolimits}
\def\colim{\mathop{\mathrm{colim}}\nolimits}
\def\Spec{\mathop{\mathrm{Spec}}}
\def\Hom{\mathop{\mathrm{Hom}}\nolimits}
\def\Ext{\mathop{\mathrm{Ext}}\nolimits}
\def\SheafHom{\mathop{\mathcal{H}\!\mathit{om}}\nolimits}
\def\SheafExt{\mathop{\mathcal{E}\!\mathit{xt}}\nolimits}
\def\Sch{\mathit{Sch}}
\def\Mor{\mathop{\mathrm{Mor}}\nolimits}
\def\Ob{\mathop{\mathrm{Ob}}\nolimits}
\def\Sh{\mathop{\mathit{Sh}}\nolimits}
\def\NL{\mathop{N\!L}\nolimits}
\def\CH{\mathop{\mathrm{CH}}\nolimits}
\def\proetale{{pro\text{-}\acute{e}tale}}
\def\etale{{\acute{e}tale}}
\def\QCoh{\mathit{QCoh}}
\def\Ker{\mathop{\mathrm{Ker}}}
\def\Im{\mathop{\mathrm{Im}}}
\def\Coker{\mathop{\mathrm{Coker}}}
\def\Coim{\mathop{\mathrm{Coim}}}

% Boxtimes
%
\DeclareMathSymbol{\boxtimes}{\mathbin}{AMSa}{"02}

%
% Macros for moduli stacks/spaces
%
\def\QCohstack{\mathcal{QC}\!\mathit{oh}}
\def\Cohstack{\mathcal{C}\!\mathit{oh}}
\def\Spacesstack{\mathcal{S}\!\mathit{paces}}
\def\Quotfunctor{\mathrm{Quot}}
\def\Hilbfunctor{\mathrm{Hilb}}
\def\Curvesstack{\mathcal{C}\!\mathit{urves}}
\def\Polarizedstack{\mathcal{P}\!\mathit{olarized}}
\def\Complexesstack{\mathcal{C}\!\mathit{omplexes}}
% \Pic is the operator that assigns to X its picard group, usage \Pic(X)
% \Picardstack_{X/B} denotes the Picard stack of X over B
% \Picardfunctor_{X/B} denotes the Picard functor of X over B
\def\Pic{\mathop{\mathrm{Pic}}\nolimits}
\def\Picardstack{\mathcal{P}\!\mathit{ic}}
\def\Picardfunctor{\mathrm{Pic}}
\def\Deformationcategory{\mathcal{D}\!\mathit{ef}}
\usepackage{cleveref}
\crefname{originallemma}{original lemma}{original lemmas}
\Crefname{originallemma}{Original lemma}{Original lemmas}
\makeatletter
\newlabel{lemma-properties-diagonal}{{03HK}{source}{}{sourcepointer-1}{}}
\newlabel{lemma-properties-diagonal@cref}{{[originallemma][0][]03HK}{[1][0][]source}{}{}{}}
\makeatother
\crefname{originalsection}{original section}{original sections}
\Crefname{originalsection}{Original section}{Original sections}
\makeatletter
\newlabel{section-representable}{{03HA}{source}{}{sourcepointer-2}{}}
\newlabel{section-representable@cref}{{[originalsection][0][]03HA}{[1][0][]source}{}{}{}}
\makeatother
\title{The comparison map between the two fibre products is a base change of the diagonal}
\author{The Stacks Project Authors}
\date{stacks-project:spaces-morphisms.tex:local-snapshot-20261012; GNU FDL 1.2 or later}
\begin{document}
\maketitle
\setcounter{subsection}{1}
\noindent\textit{Editorial note.} All delivered mathematical text is the author's own, copied from the pinned Stacks Project chapter file spaces-morphisms.tex; no statement or proof step was rewritten, reordered or removed.\par\medskip
\noindent\textit{Editorial note.} The statement begins with a slogan block (part of the counted statement span) and displays four numbered assertions; the proof identifies the base change structure once and derives all four.\par\medskip
\noindent\textit{Editorial note.} The two references of the counted spans are recorded as non-load-bearing navigation pointers; the argument that consumes them stays inside the counted proof span.\par\medskip
\noindent\textit{Editorial note.} Preamble: the three native excerpt spans of sources/stacks-spaces-morphisms/preamble.tex (lines 8-13, 14-18 and 154-217). The chapter file itself carries no preamble text.\par\medskip
\noindent\textit{Allowed standard prerequisites (editorial).}\begin{itemize}
\item The properties of the diagonal morphism of a morphism of algebraic spaces. Lemma lemma-properties-diagonal (381-394), used at proof lines 533-534. Supplies representability, local finite typeness, monomorphismhood and separatedness of the base change.
\item The square exhibiting i as a base change of the diagonal. used at proof lines 525-532. Is the categorical computation at the heart of the proof.
\item The translation between properties of representable morphisms and of their scheme base changes. Section section-representable, used at proof lines 534-535. Converts the properties of the base change into the four assertions of the statement.
\item Local separatedness, separatedness and quasi-separatedness of Z -> T. used at proof lines 518-520. Are the hypotheses under which the immersion, closed immersion and quasi-compact assertions hold.
\end{itemize}\medskip
\ifdefined\hypertarget\hypertarget{sourcepointer-1}{}\fi
\noindent\textit{Source reference.} Original lemma 03HK: focused-seeds/AGP/sources/stacks-snapshot-20261011/tags.txt, source lines 4570--4570. Navigation pointer into the Stacks Project; the standard fact it supplies is applied inside the counted proof span and the argument resting on it stays there.\par\medskip
\ifdefined\hypertarget\hypertarget{sourcepointer-2}{}\fi
\noindent\textit{Source reference.} Original section 03HA: focused-seeds/AGP/sources/stacks-snapshot-20261011/tags.txt, source lines 4559--4559. Navigation pointer into the Stacks Project; the standard fact it supplies is applied inside the counted proof span and the argument resting on it stays there.\par\medskip
\begin{lemma}
\label{lemma-fibre-product-after-map}
\begin{slogan}
The top arrow of a ``magic diagram'' of algebraic spaces has nice
immersion-like properties, and under separatedness hypotheses
these get stronger.
\end{slogan}
Let $S$ be a scheme. Let $f : X \to Z$, $g : Y \to Z$ and $Z \to T$
be morphisms of algebraic spaces over $S$. Consider the induced morphism
$i : X \times_Z Y \to X \times_T Y$. Then
\begin{enumerate}
\item $i$ is representable, locally of finite type, locally quasi-finite,
separated and a monomorphism,
\item if $Z \to T$ is locally separated, then $i$ is an immersion,
\item if $Z \to T$ is separated, then $i$ is a closed immersion, and
\item if $Z \to T$ is quasi-separated, then $i$ is quasi-compact.
\end{enumerate}
\end{lemma}

\begin{proof}
By general category theory the following diagram
$$
\xymatrix{
X \times_Z Y \ar[r]_i \ar[d] & X \times_T Y \ar[d] \\
Z \ar[r]^-{\Delta_{Z/T}} \ar[r] & Z \times_T Z
}
$$
is a fibre product diagram. Hence $i$ is the base change of the
diagonal morphism $\Delta_{Z/T}$. Thus the lemma follows
from Lemma \ref{lemma-properties-diagonal}, and the material in
Section \ref{section-representable}.
\end{proof}

\end{document}
